\chapter{Спектральный поток сфалерона и класс пятнадцать}
\label{ch:v6-spectral-transition-sphaleron-spectral-flow-gate}

\section{Вопрос о настоящем операторном мосте}

Предыдущая глава установила, что электрослабый сфалерон реализует
конечную критическую конфигурацию с $H(0)=0$, но не является устойчивой
частицей. Настоящий гейт проверяет более узкое и содержательное
утверждение: совпадает ли фермионный спектральный поток вдоль перехода
между соседними gauge-вакуумами с доказанным тёплицевым классом
$15\in KO_6$.

Необходимо различать три целых числа:
\begin{enumerate}
 \item изменение числа Черна--Саймонса $\Delta N_{\rm CS}$;
 \item число ориентированных пересечений нуля хирального гамильтониана;
 \item ранг коэффициентного проектора $q_0$ тёплицева цикла.
\end{enumerate}
Равенство двух чисел без явной карты операторов не считается мостом.

\section{Стандартный поток вдоль сфалеронной петли}

Рассмотрим путь gauge--Higgs конфигураций из вакуума с
$N_{\rm CS}=n$ в соседний вакуум с $N_{\rm CS}=n+1$. Сфалерон лежит на
вершине минимаксного барьера и имеет
\begin{equation}
 N_{\rm CS}=n+\frac12
 \qquad (\operatorname{mod}\mathbb Z).
\end{equation}
Для одного левого $SU(2)$-дублета ориентированный спектральный поток
равен единице при $\Delta N_{\rm CS}=1$; в середине симметричного пути
возникает нулевая мода. Это стандартная связь сфалерона, спектрального
потока и хиральной аномалии
\cite{sphaleron-flow-klinkhamer-rupp2003,
sphaleron-flow-klinkhamer-lee2001}.

В одном поколении Стандартной модели имеются четыре копии левого
слабого дублета:
\begin{equation}
 Q_L^r,\quad Q_L^g,\quad Q_L^b,\quad L_L.
\end{equation}
Поэтому при единичном переходе
\begin{equation}
 \operatorname{SF}_{\rm chiral}^{(1\,gen)}=3+1=4.
 \label{eq:v6-sphaleron-flow-four}
\end{equation}
Правые синглеты $u_R,d_R,e_R$ не добавляют независимых слабых
дублетных пересечений.

\section{Аномальные заряды}

Каждое из трёх кварковых пересечений несёт барионный вес $1/3$, а
лептонное пересечение --- лептонный вес $1$. Следовательно,
\begin{equation}
 \Delta B=1,\qquad
 \Delta L=1,\qquad
 \Delta(B-L)=0
 \label{eq:v6-sphaleron-anomaly-one-generation}
\end{equation}
на одно поколение и единичное $\Delta N_{\rm CS}$. Для трёх поколений
получаются двенадцать дублетных пересечений, но
\begin{equation}
 \Delta B=\Delta L=3,
\end{equation}
а не класс пятнадцать. Это обычное аномальное правило, восходящее к
инстантонному и сфалеронному анализу
\cite{sphaleron-flow-thooft1976}.

\section{Почему это не разложение \texorpdfstring{$15=12+3$}{15=12+3}}

Конечный носитель проекта имеет вид
\begin{equation}
 H_{15}=Q_L\oplus L_L\oplus u_R\oplus d_R\oplus e_R,
 \qquad 15=6+2+3+3+1,
\end{equation}
и заряженный граф сохраняет компоненты
\begin{equation}
 H_q\oplus H_\ell,
 \qquad \dim H_q=12,
 \qquad \dim H_\ell=3.
\end{equation}
Сфалеронный поток ограничивается иначе:
\begin{equation}
 \operatorname{SF}(H_q)=3,
 \qquad
 \operatorname{SF}(H_\ell)=1.
 \label{eq:v6-sphaleron-flow-three-plus-one}
\end{equation}
Он считает цветовые копии левого кваркового дублета и один левый
лептонный дублет, а не комплексную размерность полного хирального плюс
правого пакета.

Число $12$ для трёх физических поколений является
$3\times(3+1)$, тогда как проектный ранг $12$ есть
$6+3+3$ внутри \emph{одного} поколения. Совпадение символа ``12'' имеет
разные индексы происхождения и не задаёт канонического изоморфизма.

\begin{theorem}[Сфалеронный поток не равен классу пятнадцать]
Стандартный электрослабый переход с $\Delta N_{\rm CS}=1$ имеет четыре
хиральных дублетных пересечения на поколение и двенадцать на три
поколения. Доказанный проектный класс пятнадцать является рангом
коэффициентной опоры одного поколения и раскладывается как $12+3$ по
полному заряженному графу. Поэтому стандартный сфалеронный спектральный
поток не реализует этот класс без дополнительного операторного
произведения, явно вставляющего $q_0$.
\end{theorem}

\section{Real/KO6-удвоение}

Ориентированная хиральная ветвь имеет ненулевой поток. Сопряжённая ветвь
с обратной ориентацией имеет противоположный поток, так что обычная сумма
полной Real-пары равна нулю:
\begin{equation}
 (+4)+(-4)=0
\end{equation}
на поколение. Это согласуется с общей структурой KO6-пары, но не
отождествляет её с парой тёплицевых индексов $(-15,+15)$.

Теорема Атьи--Патоди--Зингера связывает спектральный поток семейства
самосопряжённых дираковских операторов с индексом оператора на цилиндре
\cite{sphaleron-flow-aps1975}. Но для применения к проекту требуется
явное семейство, коэффициентный модуль и граничная карта. Одной
аддитивности индекса недостаточно.

\section{Что всё-таки подтверждено}

Сфалерон даёт физически известный пример именно той общей схемы, ради
которой проект отказался от буквальной неразрывной верёвки:
\begin{equation}
 \text{путь конфигураций}
 \longrightarrow
 \text{нулевая мода}
 \longrightarrow
 \text{ориентированный спектральный поток}.
\end{equation}
Следовательно, язык спектрального перехода не является пустой метафорой.
Но конкретное число и носитель стандартного электрослабого потока не
совпадают с проектным $q_0$.

\begin{nogocriterion}
Нельзя объявлять сфалерон механизмом рождения класса пятнадцать по
совпадению слов ``нулевая мода'', ``индекс'' и ``спектральный поток''.
Требуется явная фредгольмова или каспаровская карта, переводящая единичный
сфалеронный класс в коэффициентный Real-класс $q_0$ и одновременно
воспроизводящая хиральные аномальные веса.
\end{nogocriterion}

\section{Следующий различающий тест}

Следующий гейт
\path{version6_spectral_transition_anomaly_to_toeplitz_product_map_gate}
должен проверить условную конструкцию
\begin{equation}
 [\text{sphaleron path}]\boxtimes[q_0]
 \longmapsto \pm15
\end{equation}
и ответить, является ли $q_0$ коэффициентом уже существующего
фермионного оператора или вставляется вручную. Отдельно требуется
сопоставить результирующую карту с аномальными весами $B$ и $L$.

Стоп-критерий: если произведение даёт пятнадцать только потому, что ранг
$q_0=15$ был заранее подставлен, не выводя этот проектор из сфалеронного
гамильтониана, такой результат является внешним декорированием, а не
механизмом рождения материи.

\begin{protocol}
\begin{enumerate}
 \item $\Delta N_{\rm CS}$ соседних вакуумов: \textbf{один};
 \item поток одного левого слабого дублета: \textbf{один};
 \item поток одного поколения: \textbf{$3+1=4$};
 \item поток трёх поколений: \textbf{$12$};
 \item аномальные заряды на поколение: \textbf{$\Delta B=\Delta L=1$};
 \item сохранение $B-L$: \textbf{да};
 \item равенство потока проектному рангу $15$: \textbf{нет};
 \item равенство сфалеронного $12$ проектному кварковому рангу $12$:
 \textbf{нет, разные носители};
 \item обычный поток полной Real-пары: \textbf{ноль};
 \item общий язык перехода через нуль: \textbf{подтверждён};
 \item явная карта к тёплицеву классу: \textbf{не построена};
 \item физическое замыкание: \textbf{не достигнуто}.
\end{enumerate}
\end{protocol}

Машинный сертификат сохранён в
\path{s2t_v6_spectral_transition_sphaleron_spectral_flow_gate_results.json}.

\begin{thebibliography}{9}
\bibitem{sphaleron-flow-klinkhamer-rupp2003}
F.~R.~Klinkhamer, C.~Rupp,
\emph{Sphalerons, Spectral Flow, and Anomalies},
Journal of Mathematical Physics 44 (2003), 3619--3639;
arXiv:hep-th/0304167.

\bibitem{sphaleron-flow-klinkhamer-lee2001}
F.~R.~Klinkhamer, Y.~J.~Lee,
\emph{Spectral Flow of Chiral Fermions in Nondissipative Yang--Mills
Gauge Field Backgrounds},
Physical Review D 64 (2001), 065024; arXiv:hep-th/0104096.

\bibitem{sphaleron-flow-thooft1976}
G.~'t~Hooft,
\emph{Computation of the Quantum Effects Due to a Four-Dimensional
Pseudoparticle},
Physical Review D 14 (1976), 3432--3450.

\bibitem{sphaleron-flow-aps1975}
M.~F.~Atiyah, V.~K.~Patodi, I.~M.~Singer,
\emph{Spectral Asymmetry and Riemannian Geometry I},
Mathematical Proceedings of the Cambridge Philosophical Society 77
(1975), 43--69.
\end{thebibliography}